\documentclass[a4j,11pt]{jarticle}
\usepackage{fancyheadings}
%\usepackage{here}
\usepackage[dvipdfmx]{graphicx}
\usepackage{listings}
\usepackage{xcolor}
\lstset{
  language=C,
  basicstyle=\ttfamily\small,
  numbers=left,
  numberstyle=\tiny,
  stepnumber=1,
  frame=single,
  breaklines=true,
  tabsize=2,
  commentstyle=\color{green!50!black},
  keywordstyle=\color{blue},
  stringstyle=\color{red}
}

\renewcommand{\thepage}{\small -- \arabic{page} --}
%\headrulewidth=0pt
\rhead{}
\lhead{}
\textheight 234mm
%
\usepackage{float}
\usepackage{amsmath}	% required for `\cases' (yatex added)
\usepackage{bm}
\usepackage{ascmac}	% required for `\screen' (yatex added)
\usepackage{color}

\title{プログラミング応用課題レポート}
\date{2026/5/19}

\author{36714087 角替勝信\\
34714148 Lee Chanhee\\
36719030 林遥都}
%
%
\begin{document}
\maketitle
\clearpage

\section{プログラム}
\begin{lstlisting}

\end{lstlisting}

\section{貸出状況の管理}
本の貸出状況を管理するために、BoFのデザインパターンのStateパターンを用いている。
BookStateインターフェースを作成し、AvailableState, LoanedStateクラスで継承している。
この時、Java 17の新機能であるSealed Classesを使用し、他のクラスから継承できないように工夫をしている。

このBookStateをBookクラスから管理している。
Bookクラスでは、現在の状況に応じて自動的に貸し借り、また現在の状態を得れるようにしている。

さらに、LibraryServiceクラスでは、引数で指定されたidの本の状態をLibraryDAOクラスから取得し、その情報を用いてBookクラスに貸し借りが可能か確認し、貸し借りが可能であれば変更後の状態をデータベースに書き込むようにしている。

\end{document}
